\documentclass{beamer}

\usetheme{Berlin}
\usecolortheme{Orchid}
\useoutertheme{miniframes}
\setbeamertemplate{navigation symbols}{}

\usepackage[utf8]{inputenc}
\usepackage[T1]{fontenc}
\usepackage{amsmath}
\usepackage{amssymb}
\usepackage{booktabs}
\usepackage{graphicx}
\graphicspath{{../images/}}

\usepackage{xcolor}
\usepackage{listings}
\usepackage{hyperref}
\usepackage{tikz}
\usetikzlibrary{arrows.meta,positioning,shapes.geometric,calc}

\lstset{
  language=Java,
  basicstyle=\ttfamily\small,
  keywordstyle=\color{blue},
  commentstyle=\color{gray},
  stringstyle=\color{teal},
  showstringspaces=false,
  columns=fullflexible,
  breaklines=true
}

% Per-slide source line (references shown on the slide itself).
\newcommand{\slidesources}[1]{%
  \vfill
  {\tiny\color{gray}\raggedright\sloppy\parbox{\linewidth}{Sources: #1}\par}%
}

% Unit-wise source strings for this subject (used by slides/notes).
% Path is relative to the lecture `latex/` folder (the compilation CWD).
% OOPS with Java (CSEG1044) — unit-wise sources
%
% These are short, student-facing reference strings intended to be used with:
%   - \slidesources{...}  (in beamer slides)
%   - \notesources{...}   (in lecture notes)
%
% Style inspiration: other subjects in My-Subjects/ use semicolon-separated sources
% and include an "accessed" date for web documentation.

\newcommand{\OOPSUnitISources}{Schildt, \textit{Java: A Beginner's Guide} (10e), McGraw-Hill (Java fundamentals + OOP basics).; Oracle Java Tutorials: Getting Started + Language Basics + Classes and Objects (accessed \today).; Java SE API docs (e.g., \texttt{String}, \texttt{Scanner}) (accessed \today).}

\newcommand{\OOPSUnitIISources}{Schildt, \textit{Java: The Complete Reference} (12e), McGraw-Hill (classes, inheritance, packages, interfaces).; Bloch, \textit{Effective Java} (3e), Addison-Wesley, 2018 (design choices; interfaces vs abstract classes).; Oracle Java Tutorials: Classes and Objects + Interfaces + Packages (accessed \today).; Java SE API docs (e.g., \texttt{Comparable}, \texttt{Runnable}) (accessed \today).}

\newcommand{\OOPSUnitIIISources}{Schildt, \textit{Java: The Complete Reference} (12e) (nested/inner classes, exceptions, threads, I/O).; Oracle Java Tutorials: Nested Classes + Exceptions + Concurrency + Basic I/O (accessed \today).; Java SE API docs (e.g., \texttt{Thread}, \texttt{AutoCloseable}) (accessed \today).}

\newcommand{\OOPSUnitIVSources}{Schildt, \textit{Java: The Complete Reference} (12e) (generics, lambdas, Swing, JDBC).; Bloch, \textit{Effective Java} (3e) (generics + lambdas).; Oracle Java Tutorials: Generics + Lambda Expressions + Swing + JDBC Basics (accessed \today).; Java SE API docs (e.g., \texttt{java.util.function}, \texttt{javax.swing}, \texttt{java.sql}) (accessed \today).}

\newcommand{\OOPSUnitVSources}{Schildt, \textit{Java: The Complete Reference} (12e) (Collections Framework + wrapper classes).; Oracle Java docs: Collections Framework overview (accessed \today).; Java SE API docs (\texttt{java.util} collections; wrapper classes in \texttt{java.lang}) (accessed \today).}

\newcommand{\OOPSUnitVISources}{Git SCM: \textit{Pro Git} (2e) (version control workflow), accessed \today.; GitHub Docs (repositories, issues, pull requests), accessed \today.; Oracle Java Tutorials: Swing + JDBC (accessed \today).; JUnit 5 User Guide (accessed \today).; Apache Maven Project: Getting Started (accessed \today).; Gradle User Manual: Getting Started (accessed \today).; UML class diagrams: UML 2.x overview/spec (accessed \today).}

\newcommand{\OOPSMiscWhyInterfacesSources}{Schildt, \textit{Java: The Complete Reference} (12e) (interfaces).; Bloch, \textit{Effective Java} (3e) (prefer interfaces where appropriate).; Oracle Java Tutorials: Interfaces (accessed \today).; Java SE API docs (e.g., \texttt{Comparable}, \texttt{Runnable}) (accessed \today).}



\title[OOPS with Java]{Object Oriented Programming (CSEG1044)}
\subtitle{Miscellaneous: Why Do We Need Interfaces?}
\author{Tofik Ali}
\institute{School of Computer Science, UPES Dehradun}
\date{\today}

\begin{document}

\begin{frame}[fragile]
  \titlepage
  \vspace{0.5em}
  \slidesources{\OOPSMiscWhyInterfacesSources}
\end{frame}

\begin{frame}{Quick Links}
  \centering
  \hyperlink{sec:overview}{\beamerbutton{Overview}}\hspace{0.6em}
  \hyperlink{sec:motivation}{\beamerbutton{Motivation}}\hspace{0.6em}
  \hyperlink{sec:usecases}{\beamerbutton{Use Cases}}\hspace{0.6em}
  \hyperlink{sec:summary}{\beamerbutton{Summary}}
\end{frame}

\begin{frame}{Agenda}
  \tableofcontents
\end{frame}

\section{Overview}
\label{sec:overview}

\begin{frame}{Learning Outcomes}
  \begin{itemize}[<+->]
    \item Explain what an interface represents (a contract for behavior).
    \item Explain why interfaces reduce coupling and improve design flexibility.
    \item Apply interfaces in common scenarios: strategy, testing, and data access.
    \item Decide when to use an interface vs an abstract class (rule-of-thumb).
  \end{itemize}
\end{frame}

\begin{frame}{The one-line answer}
  \begin{block}{Why do we need interfaces?}
    So our code depends on \textbf{contracts} (behavior) rather than \textbf{concrete classes} (implementation).
  \end{block}

  \vspace{0.5em}
  \begin{itemize}[<+->]
    \item Easier to change (swap implementations without rewriting business logic)
    \item Easier to extend (add new implementations without touching existing code)
    \item Easier to test (replace real dependencies with fakes/mocks)
  \end{itemize}
\end{frame}

\section{Motivation (Tight Coupling)}
\label{sec:motivation}

\begin{frame}[fragile]{Without interfaces: a rigid design}
\begin{lstlisting}
class CardPayment {
    void pay(int amount) { /* call bank */ }
}

class CheckoutService {
    void checkout(int amount) {
        CardPayment payment = new CardPayment(); // fixed choice
        payment.pay(amount);
    }
}
\end{lstlisting}
\end{frame}

\begin{frame}{Symptoms of tight coupling}
  \begin{itemize}[<+->]
    \item New requirement (UPI/Wallet) forces edits in \texttt{CheckoutService}
    \item Testing is painful (real bank calls, slow and flaky)
    \item High-level business code depends on low-level implementation details
  \end{itemize}
\end{frame}

\section{Contract (Interface)}
\label{sec:contract}

\begin{frame}[fragile]{Introduce a contract}
\begin{lstlisting}
interface PaymentMethod {
    void pay(int amount);
}

class CardPayment implements PaymentMethod {
    public void pay(int amount) { /* call bank */ }
}

class UpiPayment implements PaymentMethod {
    public void pay(int amount) { /* call UPI */ }
}
\end{lstlisting}
\end{frame}

\begin{frame}[fragile]{Business logic depends on the contract}
\begin{lstlisting}
class CheckoutService {
    private final PaymentMethod payment;

    CheckoutService(PaymentMethod payment) {
        this.payment = payment;
    }

    void checkout(int amount) {
        payment.pay(amount);
    }
}
\end{lstlisting}
\end{frame}

\begin{frame}{What we gained}
  \begin{itemize}[<+->]
    \item \textbf{Flexibility:} choose implementation at runtime (Card/UPI/Wallet)
    \item \textbf{Extensibility:} add a new payment type without changing checkout
    \item \textbf{Testability:} pass a fake implementation in tests
  \end{itemize}

  \vspace{0.5em}
  \begin{center}
    \texttt{CheckoutService} $\rightarrow$ \texttt{PaymentMethod} $\leftarrow$ \texttt{CardPayment / UpiPayment}
  \end{center}
\end{frame}

\section{Use Cases}
\label{sec:usecases}

\begin{frame}[fragile]{Use case 1: plug-in behavior (Strategy)}
\begin{lstlisting}
interface DiscountPolicy { int discount(int total); }
class NoDiscount implements DiscountPolicy {
    public int discount(int total) { return 0; }
}
class StudentDiscount implements DiscountPolicy {
    public int discount(int total) { return total / 10; } // 10%
}
class Billing {
    private final DiscountPolicy policy;
    Billing(DiscountPolicy policy) { this.policy = policy; }
    int finalAmount(int total) { return total - policy.discount(total); }
}
\end{lstlisting}
\end{frame}

\begin{frame}[fragile]{Use case 2: unit testing with a fake}
\begin{lstlisting}
interface Notifier { void send(String to, String msg); }

class FakeNotifier implements Notifier {
    String lastTo, lastMsg;
    public void send(String to, String msg) {
        lastTo = to; lastMsg = msg;
    }
}
\end{lstlisting}
\begin{itemize}[<+->]
  \item Real code uses \texttt{EmailNotifier}; tests use \texttt{FakeNotifier}
  \item No real network calls; tests become fast and reliable
\end{itemize}
\end{frame}

\begin{frame}[fragile]{Use case 3: decouple storage (Repository)}
\begin{lstlisting}
class User { String id; String name; }

interface UserRepository {
    User findById(String id);
    void save(User u);
}

class InMemoryUserRepository implements UserRepository { /* for tests */ }
class SqlUserRepository implements UserRepository { /* for production */ }
\end{lstlisting}
\end{frame}

\section{Java Examples}
\label{sec:java}

\begin{frame}{Interfaces you already use in Java}
  \begin{itemize}[<+->]
    \item Collections: \texttt{List}, \texttt{Set}, \texttt{Map}, \texttt{Iterable}, \texttt{Iterator}
    \item Ordering: \texttt{Comparable}, \texttt{Comparator}
    \item Concurrency: \texttt{Runnable}, \texttt{Callable}
    \item Resource handling: \texttt{AutoCloseable} (try-with-resources)
  \end{itemize}
\end{frame}

\begin{frame}{Why Java loves interfaces}
  \begin{itemize}[<+->]
    \item A class can implement \textbf{multiple interfaces} (multiple roles/capabilities)
    \item Frameworks use interfaces for callbacks (``I will call your code'')
    \item Standard libraries are designed around interface types (program to the abstraction)
  \end{itemize}
\end{frame}

\section{Summary}
\label{sec:summary}

\begin{frame}{Interface vs abstract class (simple rule)}
  \begin{itemize}[<+->]
    \item Use an \textbf{interface} for a capability/role with many unrelated implementations.
    \item Use an \textbf{abstract class} when you need shared state + shared code for closely related classes.
    \item Interfaces can have \texttt{default} methods; use them carefully to keep the contract clean.
  \end{itemize}
\end{frame}

\begin{frame}{Summary}
  \begin{itemize}[<+->]
    \item Interfaces reduce coupling by separating \textit{what} from \textit{how}.
    \item They enable plug-in designs (strategy), test doubles (fakes/mocks), and clean architectures (repositories).
    \item Many Java APIs are built around interface types.
  \end{itemize}
\end{frame}

\begin{frame}{Quick practice}
  \begin{enumerate}[<+->]
    \item Design an \texttt{Exporter} interface with \texttt{CsvExporter} and \texttt{JsonExporter}.
    \item Write a tiny test that uses a \texttt{FakeNotifier} and checks the last message sent.
    \item Refactor any earlier program: replace direct concrete usage with an interface + two implementations.
  \end{enumerate}
\end{frame}

\end{document}
